%% kapa: exam=1 round=19 year=2008 period=1 form=A subject=civil_law subject_ko=민법 seq=28 number=28
%% source: Q-Net 감정평가사 기출문제 (cst00309) artlSeq=1001962 file=extracted/1008/1001962/2008년 제19회 감정평가사 1차시험 기출문제/1교시.hwp
%% conversion: kordoc 4.24.51 -> tools/gen1.py (auto); review=proofread-sonnet; flags=none
\begin{question}{28}
점유에 관한 설명으로 틀린 것은? (다툼이 있는 경우에는 판례에 의함)
\begin{choices}
  \item 직접점유자가 임의로 점유를 타인에게 양도한 경우에는 그 점유이전이 간접점유자의 의사에 반한다고 하더라도 간접점유자의 점유가 침탈된 경우에 해당하지 않는다.
  \item 타주점유가 자주점유로 되기 위해서는 점유자가 소유자에 대하여 소유의 의사가 있는 것을 표시하거나 새로운 권원에 의하여 다시 소유의 의사로써 점유를 시작하여야 한다.
  \item 점유자가 취득시효기간 경과 후, 등기부상의 권리자에게 매수의사를 표시하는 것은 특별한 사정이 없는 한 시효이익의 포기로 볼 수 없다.
  \item 지상건물과 함께 그 대지를 매수한 자가 착오로 인접 토지의 일부를 그가 매수․취득한 대지에 속하는 것으로 믿고 인도받아 점유하고 있는 경우는 자주점유로 볼 수 있다.
  \item 점유자의 점유가 소유의 의사 있는 자주점유인지 아니면 소유의 의사 없는 타주점유인지의 여부는 점유자의 내심의 의사에 의하여 결정된다.
\end{choices}
\end{question}
\answer{5}
